A la fin de ce chapitre, je saurai :
\begin{competences}
    \point Citer et exploiter la relation entre l’enthalpie libre de     réaction et les potentiels des couples mis en jeu dans une réaction d’oxydoréduction.
    \point Déterminer l’enthalpie libre standard d’une réaction d’oxydoréduction à partir des potentiels standard des couples.
    \point Déterminer la valeur du potentiel standard d’un couple d’oxydoréduction à partir de données thermodynamiques.
    \point Décrire le montage à trois électrodes permettant de tracer des courbes courant-potentiel.
    \point Relier vitesse de réaction électrochimique et intensité du courant.
    \point Identifier le caractère lent ou rapide d’un système à partir des courbes courant-potentiel.
    \point Identifier les espèces électroactive pouvant donner lieu à une limitation en courant par diffusion.
    \point Identifier des paliers de diffusion limite sur des relevés expérimentaux.
    \point Relier, à l’aide de la loi de Fick, l'intensité du courant de diffusion limite à la concentration du réactif et à la surface immergée de l'électrode.
    \point Tracer l’allure de courbes courant-potentiel de branches d’oxydation ou de réduction à partir de données fournies, de potentiels standard, concentrations et surpotentiels.
    \point \faSignLanguage Tracer et exploiter des courbes courant-potentiel.
    \point Établir l’inégalité reliant la variation d’enthalpie libre et le travail électrique.
    \point Relier la tension à vide d’une pile et l’enthalpie libre de la réaction modélisant son fonctionnement.
    \point Déterminer la capacité électrique d’une pile.
    \point Exploiter les courbes courant-potentiel pour rendre compte du fonctionnement d’une pile électrochimique et tracer sa caractéristique.
    \point Citer les paramètres influençant la résistance interne d’une pile électrochimique.
    \point Exploiter les courbes courant-potentiel pour rendre compte du fonctionnement d’un électrolyseur et prévoir la valeur de la tension minimale à imposer.
    \point Exploiter les courbes courant-potentiel pour justifier les contraintes (purification de la solution électrolytique, choix des électrodes) dans la recharge d’un accumulateur.
    \point Déterminer la masse de produit formé pour une durée et des conditions données d’électrolyse.
    \point Déterminer un rendement faradique à partir d’informations fournies concernant le dispositif étudié.
    \point \faSignLanguage Étudier le fonctionnement d’une pile ou d’un électrolyseur pour effectuer des bilans de matière et des bilans électriques.
    \point Positionner un potentiel de corrosion sur un tracé de courbes courant-potentiel.
    \point Interpréter le phénomène de corrosion uniforme d’un métal ou de deux métaux en contact en utilisant des courbes courant-potentiel ou d’autres données expérimentales, thermodynamiques et cinétiques.
    \point Déterminer une vitesse de corrosion.
    \point Citer des facteurs favorisant la corrosion.
    \point Exploiter des tracés de courbes courant-potentiel pour expliquer qualitativement :
    \begin{itemize}
        \item la qualité de la protection par un revêtement métallique ;
        \item le fonctionnement d’une anode sacrificielle.
    \end{itemize}
    \point Interpréter le phénomène de passivation sur une courbe courant-potentiel.
    \point \faSignLanguage Mettre en évidence le phénomène de corrosion et les facteurs l’influençant.
\end{competences}